%%
% 第四章 P450 专属模型训练与公平评估
% 核心实证章。数据集 + 公平对比是本论文的核心贡献。
%%

\chapter{P450 专属模型训练与公平评估}
\label{cha:training}

本章以第三章构建的 P450 专属基准数据集为基础，开展两类实证实验：其一，在领域专属数据上从头训练与 EZSpecificity 同架构的原架构模型（以下简称原架构模型或 M0），评估领域专属训练在 P450 家族上的性能上界；其二，以 EZSpecificity 公开发布的预训练 checkpoint 在本研究测试集上进行外部评估，并通过泄漏控制的酶层面过滤方法与本研究基线模型进行公平对比。本章结果是全文的核心实证，直接回答第一章提出的研究问题：通用 SOTA 模型 EZSpecificity 在 ESIBank 中覆盖稀缺的 P450 家族上，其 对未见酶的泛化能力是否足以支撑其"通用性"声明；若不足，则领域专属训练能在多大程度上突破这一瓶颈。

\section{实验配置训练流程与指标计算}
\label{sec:ch4-config}

\textbf{硬件与训练框架}。模型训练在 4 块 NVIDIA RTX 5090 显卡上采用分布式数据并行（Distributed Data Parallel，DDP）策略进行。训练框架基于 PyTorch Lightning\cite{Lightning2019}，模型与数据管线完全复用 EZSpecificity 官方实现\cite{EZSpec2025}，仅将数据加载器替换为第三章特征管线生成的 LMDB 与二进制存储以充分发挥 DDP 吞吐。

\textbf{训练超参数}。单卡批大小为 88、有效全局批大小为 352（88 × 4 卡）；优化器采用 AdamW\cite{AdamW2019}，一次训练至多 200 个 epoch，并以验证集 AUPR 为监控指标、耐心（patience）15 epoch 的早停策略；学习率调度采用 ReduceLROnPlateau，在验证指标连续多轮停滞后按固定比例衰减；混合精度训练以节省显存。损失函数为带类别权重的二元交叉熵（binary cross-entropy with logits），以缓解第三章所述 $1:9.9$ 的正负样本不均衡。

\textbf{评估指标}。报告两个指标。AUC-ROC（area under the receiver operating characteristic curve，受试者工作特征曲线下面积）是成对排序意义下的稳健指标，但 Saito 与 Rehmsmeier 指出其在极不平衡数据上的绝对值可能偏乐观、单独使用易产生误导\cite{Saito2015}；AUPR（area under the precision--recall curve，精确率--召回率曲线下面积）则在极度不平衡数据下对稀有正样本的打分分布更敏感，其随机基线即正样本占比，绝对数值往往远低于 AUC。由于本研究测试集正负比约为 $1:10.2$，AUC 与 AUPR 并列报告可避免单一指标掩盖模型在稀有类上的真实表现。

\begin{table}[htbp]
  \centering
  \caption{本研究基线模型训练配置}
  \label{tab:ch4-config}
  \begin{tabular}{ll}
    \toprule
    配置项 & 取值 \\
    \midrule
    硬件环境     & 4 $\times$ NVIDIA RTX 5090，DDP \\
    训练框架     & PyTorch Lightning \\
    批大小       & 88 per GPU，有效 352 \\
    优化器       & AdamW \\
    最大 epoch   & 200 \\
    早停策略     & 监控 val AUPR，patience = 15 \\
    学习率调度   & ReduceLROnPlateau \\
    损失函数     & 带类别权重的 BCEWithLogits \\
    混合精度     & 启用 \\
    切分         & random split fold 0 \\
    \bottomrule
  \end{tabular}
\end{table}

\section{原架构基线模型训练结果}
\label{sec:ch4-baseline}

原架构模型 M0 完全复用 EZSpecificity\cite{EZSpec2025} 的三通道并行架构：ESM-2\cite{ESM2023} 提供的酶残基级序列嵌入构成序列通道；GROVER\cite{GROVER2020} 与 Morgan/ECFP\cite{ECFP2010} 共同构成底物的分子通道；对接复合物经 SE(3)-等变图神经网络编码后构成结构通道。三通道经双层交叉注意力融合后经多层感知机输出二分类 logits。

模型在 4 块 RTX 5090 上完成 57 个 epoch 即因早停策略终止，总训练时长约 56.8 分钟；验证集最佳 checkpoint 出现在第 43 轮，对应 Val AUC $= 0.9250$。以该 checkpoint 在随机切分第一折测试集（10{,}999 条样本、984 正 / 10{,}015 负）上评估，得到 \textbf{Test AUC $= 0.9302$、Test AUPR $= 0.6749$}（见表~\ref{tab:ch4-baseline}）。Test AUC 数值上与 Val AUC 差距极小，表明原架构模型在本数据集上的训练并未出现明显的过拟合或分布偏移问题。

\begin{figure}[htbp]
  \centering
  \fbox{\parbox[c][3.5cm][c]{0.85\textwidth}{\centering 图 4-1 占位：原架构基线模型训练动态曲线\\[0.5em] （train/val loss、val AUC、val AUPR 随 epoch 变化）\\[0.5em] 待插入：\texttt{image/fig4-1-training-dynamics.png}}}
  \caption{原架构基线模型（M0）训练动态曲线}
  \label{fig:ch4-training}
\end{figure}

\begin{table}[htbp]
  \centering
  \caption{原架构基线模型（M0）在 random split fold 0 测试集上的结果}
  \label{tab:ch4-baseline}
  \begin{tabular}{lc}
    \toprule
    指标 / 配置 & 数值 \\
    \midrule
    最佳 checkpoint   & ep43 \\
    Val AUC           & 0.9250 \\
    Test AUC          & \textbf{0.9302} \\
    Test AUPR         & \textbf{0.6749} \\
    测试样本          & 10{,}999 \\
    正样本            & 984 \\
    负样本            & 10{,}015 \\
    正负比            & $1:10.2$ \\
    \bottomrule
  \end{tabular}
\end{table}

\section{参考论文模型的外部测试表现}
\label{sec:ch4-external}

为评估 EZSpecificity 作为 SOTA 通用模型在 P450 家族上的 对未见酶的泛化能力，本研究直接加载 EZSpecificity 论文公开发布的预训练 checkpoint\cite{EZSpec2025}，以其原始训练超参在本研究测试集上推理。重要前提：EZSpecificity 原始训练集 ESIBank 已包含 389 个 P450 酶（详见第三章），这些酶与本研究测试集中的部分酶存在重合，因此需要同时报告未过滤与过滤后两种设置下的结果（详见第 \ref{sec:ch4-fair} 节）。

在未过滤的完整测试集（10{,}999 条样本）上，参考论文模型得到 \textbf{Test AUC $= 0.5860$、Test AUPR $= 0.1124$}；在过滤掉 389 个与原训练集重合的 P450 酶后得到的 7{,}963 条样本测试集上，模型表现进一步下降至 \textbf{Test AUC $= 0.5586$、Test AUPR $= 0.1004$}。AUPR 指标的数值接近 $1:10.2$ 正负比下纯随机预测的理论基线 $\approx 0.089$，提示参考模型在本研究的 P450 测试集上几乎没有有意义的底物--酶区分能力。

\begin{table}[htbp]
  \centering
  \caption{EZSpecificity 公开 checkpoint 在本研究测试集上的表现}
  \label{tab:ch4-paper}
  \begin{tabular}{lrcc}
    \toprule
    测试集 & 样本数 & AUC & AUPR \\
    \midrule
    未过滤     & 10{,}999 & 0.5860 & 0.1124 \\
    过滤后     & 7{,}963  & \textbf{0.5586} & \textbf{0.1004} \\
    \bottomrule
  \end{tabular}
\end{table}

\section{泄漏控制条件下的公平对比}
\label{sec:ch4-fair}

\subsection{泄漏控制的必要性}

在对同一预训练 checkpoint 做外部评估时，若测试集中的酶与训练集酶存在重合，评估将无法区分两种本质不同的模型能力：对已在训练中见过的酶做推理所反映的是"记忆效应"而非"泛化能力"；仅对训练中未见过的酶做推理得到的结果才真正度量该模型在新酶上的底物特异性预测能力。第三章已详述本研究采用 InterPro\cite{Blum2025} 功能域对 ESIBank 原始 25{,}225 条酶记录严格识别出 389 个真正的 P450，并据此从本研究测试集中剔除对应酶样本得到过滤后测试集，此举将参考论文模型在本测试集上的评估条件收敛为"完全未见过的 P450 酶"。

\subsection{三路评估与差值解读}

表~\ref{tab:ch4-fair} 汇总了参考论文模型与本研究基线模型在两种测试集上的完整评估结果，构成本论文的核心公平对比。关键观察有三。

第一，\textbf{参考论文模型的性能塌缩}。从未过滤到过滤后，参考模型 AUC 由 $0.5860$ 下降至 $0.5586$（$\Delta = -0.0274$），AUPR 由 $0.1124$ 下降至 $0.1004$（$\Delta = -0.0120$）。两项指标的同向下降印证了记忆效应在未过滤条件下对评估结果的贡献：过滤掉重合的 P450 酶后，参考模型在 P450 家族上对未见酶的真实泛化能力暴露为接近随机水平。

第二，\textbf{本研究基线模型的稳健性}。本研究基线模型 M0 在未过滤测试集上 Test AUC $= 0.9302$、AUPR $= 0.6749$；过滤后测试集上 Test AUC $= 0.9205$、AUPR $= 0.6403$。两项指标的下降幅度（$\Delta\mathrm{AUC} = -0.0097$、$\Delta\mathrm{AUPR} = -0.0346$）显著小于参考模型相应的下降幅度，且绝对数值仍远高于参考模型。这一对比表明本研究基线模型的预测能力主要来自对 P450 家族本身特异性规律的学习，而非对特定已见酶的记忆。

第三，\textbf{公平对比的核心差值}。在统一的过滤后测试集（7{,}963 条样本、715 正 / 7{,}248 负）上，本研究基线模型 $0.9205$ 较参考论文模型 $0.5586$ 高出 $\Delta\mathrm{AUC} = +0.3619$；AUPR 维度高出 $\Delta = +0.5399$。这一差距较参考模型未过滤时的 $\Delta\mathrm{AUC} = +0.3442$ 进一步放大，印证过滤协议有效剔除了记忆效应的干扰，体现的是真实的泛化能力差距。

\begin{table}[htbp]
  \centering
  \caption{泄漏控制条件下的公平对比完整结果（random split fold 0）}
  \label{tab:ch4-fair}
  \begin{tabular}{llrrr}
    \toprule
    模型 & 测试集 & 样本数 & AUC & AUPR \\
    \midrule
    参考论文模型      & 未过滤 & 10{,}999 & 0.5860 & 0.1124 \\
    参考论文模型      & 过滤后 & 7{,}963  & \textbf{0.5586} & \textbf{0.1004} \\
    本研究 M0 基线    & 未过滤 & 10{,}999 & 0.9302 & 0.6749 \\
    本研究 M0 基线    & 过滤后 & 7{,}963  & \textbf{0.9205} & \textbf{0.6403} \\
    \midrule
    \multicolumn{2}{l}{$\Delta$（本研究 $-$ 参考论文，过滤后）} & --- & $+0.3619$ & $+0.5399$ \\
    \bottomrule
  \end{tabular}
\end{table}

\begin{figure}[htbp]
  \centering
  \fbox{\parbox[c][4cm][c]{0.85\textwidth}{\centering 图 4-2 占位：公平对比柱状图\\[0.5em] 四条柱对比：参考论文模型未过滤 / 过滤后 vs 本研究 M0 未过滤 / 过滤后\\[0.5em] 左轴 AUC，右轴 AUPR\\[0.5em] 脚本待运行：\texttt{image/fig4-2-fair-comparison.png}}}
  \caption{泄漏控制外部公平对比可视化}
  \label{fig:ch4-bar}
\end{figure}

\section{结果分析与应用意义}
\label{sec:ch4-discussion}

\subsection{SOTA 通用模型在稀缺家族上的泛化瓶颈}

上述公平对比得到两条明确结论。其一，EZSpecificity 作为通用酶底物特异性预测 SOTA 模型，在其训练集 ESIBank 中覆盖稀缺（仅 389 条、占 1.5\%）的 P450 超家族上，过滤酶重合后在本研究测试集上的 AUC 仅 $0.5586$、AUPR 仅 $0.1004$，后者接近 $1:10.2$ 正负比下的随机预测基线。该结果提示：通用模型在原训练集欠覆盖的酶家族上，即使架构先进，对未见酶的泛化能力仍面临显著瓶颈。其二，以同架构在 P450 专属数据上从头训练的本研究基线模型在过滤后测试集上达到 AUC $= 0.9205$、AUPR $= 0.6403$，相对参考论文模型有 $+0.3619$ 的 AUC 差距与 $+0.5399$ 的 AUPR 差距。这一经验观察支持以下方法学结论：面向训练集中覆盖稀缺的特定酶家族，领域专属数据的微调对于突破通用模型的 未见酶的泛化瓶颈是必要的。

需要强调的是，本研究的结论严格限定为\emph{在本 P450 专属测试集与过滤协议下的观察}，不对 EZSpecificity 论文在其原始报告的卤化酶及其他六个代表性家族上的性能作任何重新评价。EZSpecificity 的通用性声明在其原始数据覆盖范围内仍然有效；本研究的贡献在于在一个数据覆盖显著稀疏的 P450 家族上，具体量化该声明的边界。

\subsection{对药学筛查流程的参考意义}

上述结论对药学领域的 P450 计算筛查实践具有直接参考意义。首先，\textbf{在将通用预训练酶--底物特异性预测模型应用于 P450 场景时，必须独立评估其在领域专属数据上的泛化能力}，不可直接依据原论文报告的整体指标推断其在 P450 上的可信度；在无法进行独立评估时，应将模型输出作为低优先级的辅助参考，而非筛查决策依据。其次，\textbf{在算力与数据可得的条件下，基于公开架构与 P450 专属数据的重训基线是一项工程上可行的工具}，其性能在本测试集上显著超越通用 checkpoint 在未见酶上的预测结果，可作为虚拟筛选早期阶段的候选 P450 底物识别工具。

作为文献讨论，以下给出三个具体药学场景以说明 P450 底物特异性预测的潜在应用方向。第一，\textbf{药物代谢早期筛选}。CYP3A4 代谢约半数临床药物，包括辛伐他汀、阿托伐他汀、咪达唑仑等。药物--药物相互作用（DDI）引发的 CYP3A4 抑制或诱导可导致严重不良反应，典型案例如西沙比利（cisapride）作为 CYP3A4 的底物与 CYP3A4 抑制剂（如伊曲康唑）联用时血药浓度升高、引发致死性 QT 延长的早期病例报告\cite{Wysowski1996} 最终促成其在 2000 年全球撤市。具备 P450 专属泛化能力的预测模型可在先导化合物阶段对分子与主要 CYP 同工酶的亲和性做快速预估。第二，\textbf{药物相互作用与遗传多态性预警}。CYP2D6 具有显著的遗传多态性\cite{Zanger2013}，相应酶活性跨 poor / intermediate / extensive / ultra-rapid 四种代谢表型变化；前药可待因依赖 CYP2D6 代谢活化为吗啡，超快代谢者服用可致中毒，而弱代谢者镇痛不足；他莫昔芬活化为 endoxifen 同样依赖 CYP2D6\cite{Stearns2003}。P450 特异性预测可辅助识别 CYP2D6 底物候选化合物以支持临床个体化用药决策。第三，\textbf{植物天然产物合成生物学}。青蒿素、紫杉醇、人参皂苷等重要天然药物的生物合成路径中多含 P450 氧化步骤（如紫杉醇合成涉及 CYP725A 家族的多个 P450 介导的连续羟化步骤\cite{Kaspera2006}），在异源表达底盘（酵母、烟草等）中挑选合适的植物 P450 酶是合成生物学改造的核心瓶颈；本研究纳入的 PlantP450DB 与 PCPD 来源大幅扩展了植物 P450 的数据覆盖，使得面向植物 P450 的特异性预测具备更充分的领域数据基础。

需要说明的是，上述应用均为基于本研究模型的潜在应用方向，具体药物研发场景仍需结合体外酶学实验与体内药理毒理验证；本论文不对临床决策作出直接断言。本研究的工程贡献在于为上述场景提供一个经严格公平评估的候选工具，以及一个可供社群复用的 P450 专属基准数据集。
